% META: {"id": "1SPE-PROBCOND-TD-001", "chapitre": "1SPE-PROBA-COND", "type_objet": "cours", "sous_type": "td_contextualise", "capacites_codes": ["C1","C2","C3","C5"], "sources_inspiration": [], "status": "generated"}
% BEGIN-VERIFY
% from sympy import Rational, simplify
% # Partie A : prevalence 2%, sensibilite 90%, specificite 95%
% P_M = Rational(2, 100)
% P_Mbar = 1 - P_M
% P_T_M = Rational(90, 100)
% P_Tbar_Mbar = Rational(95, 100)
% P_T_Mbar = 1 - P_Tbar_Mbar
% assert P_T_Mbar == Rational(1, 20)
% # P(T) par prob totales
% P_T = P_M * P_T_M + P_Mbar * P_T_Mbar
% assert P_T == Rational(2,100)*Rational(90,100) + Rational(98,100)*Rational(5,100)
% assert P_T == Rational(180, 10000) + Rational(490, 10000)
% assert P_T == Rational(670, 10000)
% assert P_T == Rational(67, 1000)
% # VPP
% P_M_T = P_M * P_T_M / P_T
% assert P_M_T == Rational(180, 10000) / Rational(67, 1000)
% assert P_M_T == Rational(18, 67)
% # Partie B : prevalence 20%
% P_M2 = Rational(20, 100)
% P_Mbar2 = 1 - P_M2
% P_T2 = P_M2 * P_T_M + P_Mbar2 * P_T_Mbar
% assert P_T2 == Rational(20,100)*Rational(90,100) + Rational(80,100)*Rational(5,100)
% assert P_T2 == Rational(1800, 10000) + Rational(400, 10000)
% assert P_T2 == Rational(2200, 10000)
% assert P_T2 == Rational(11, 50)
% P_M_T2 = P_M2 * P_T_M / P_T2
% assert P_M_T2 == Rational(9, 11)
% END-VERIFY

\section*{TD --- Depistage et valeur predictive}

\textit{Duree estimee : 50 minutes. Capacites : C1, C2, C3, C5.}

\bigskip

Un laboratoire evalue un test de depistage pour une maladie rare. Le test a une \textbf{sensibilite} de $90\,\%$ (il detecte $90\,\%$ des malades) et une \textbf{specificite} de $95\,\%$ (il identifie correctement $95\,\%$ des personnes saines).

On note $M$ l'evenement << la personne est malade >> et $T$ l'evenement << le test est positif >>.

\subsection*{Partie A --- Maladie rare (prevalence $2\,\%$)}

Dans la population generale, $2\,\%$ des personnes sont atteintes.

\begin{enumerate}
  \item Traduire les donnees de l'enonce en probabilites : $P(M)$, $P_M(T)$, $P_{\overline{M}}(T)$.

  \item Construire un arbre pondere complet a deux niveaux.

  \item Calculer $P(T)$ a l'aide de la formule des probabilites totales.

  \item Calculer la \textbf{valeur predictive positive} $P_T(M)$. Interpreter le resultat.

  \item Expliquer pourquoi, malgre la bonne qualite du test, un resultat positif n'est pas synonyme de maladie certaine.
\end{enumerate}

\subsection*{Partie B --- Maladie frequente (prevalence $20\,\%$)}

On suppose maintenant que la prevalence est de $20\,\%$.

\begin{enumerate}
  \setcounter{enumi}{5}
  \item Recalculer $P(T)$ et $P_T(M)$ avec cette nouvelle prevalence.

  \item Comparer avec la partie A et commenter l'influence de la prevalence.
\end{enumerate}

\subsection*{Partie C --- Synthese}

\begin{enumerate}
  \setcounter{enumi}{7}
  \item Resumer dans un tableau les valeurs de $P_T(M)$ pour les prevalences $2\,\%$ et $20\,\%$.

  \item Un depistage systematique de la population generale est-il pertinent pour une maladie rare ? Argumenter.
\end{enumerate}
